\cotefile{109}
\scanpage{11}
a suitable notion of a ``tame'' topological space, viewed as the natural set-up for expressing and stimulating the topological intuition of shape and form, which for so long has been squeezed \uline{into} the straitjacket of a ``General Topology'' tailored for the use of the analyst, not for geometry. This is inapt notably (because of pervasive ``wildness'' phenomena) for \uline{\inserted{the} study\deleted{ing} \inserted{of}} homotopy types of spaces of homeomorphisms, imbeddings \& the like, or for developing a comprehensive ``dévissage'' theory of stratified structures (whereas such structures \deleted{abound} \uline{crop up everywhere} notably in any question of ``moduli'' \deleted{\unsure{of}}for geometrical configurations capable of continuous variation and giving rise to degeneracy phenomena). The other continent appears \uline{at} the confluence of a multitude of streams of thought and insights, coming from topology, conformal geometry, algebraic geometry, discrete and profinite groups, algebraic groups over number fields, regular polyhedra, arithmetics, ``motives''\ldots, \uline{called together for} the exploration of the action of profinite ``absolute'' Galois groups (such as $\operatorname{Gal}_{\overline{\mathbb{Q}}/\mathbb{Q}}$) on certain (so called ``anabelian'') profinite geometrical fundamental groups (such as $\pi_{1}(P_{\overline{\mathbb{Q}}}-\{0,1,\infty\})$,

One of the leading threads comes from the combinatorial topology

(Continues without change until last paragraph)

The \inserted{two} first \deleted{two} volumes ``Pursuing Stacks'', started\textendash{} through the impetus of a correspondence in English are written in that language. They begin with a substantial common introduction \uline{of which a first part}, of personal rather than technical character, is written in French (the author's adopted language.).
% archive sheet number: 10
